\section{Modeling and validation}
\label{sec:modelacion}

All techniques share the same temporal partition: the years 2021--2024 are
used for fitting and the full year 2025 for validation. The rule is strict:
no procedure sees 2025 during fitting, including PCA and FA, whose axes are
estimated on 2021--2024 and applied to 2025 with the training weights;
otherwise information from the validation partition would leak into the
coordinates consumed by the classifiers. The only remaining leak is the
Box-Cox parameter $\lambda$, estimated on the full period in
\cref{sec:transformaciones}; it is declared because reverting it would
require rebuilding the transformed dataset.

Validation has a particularity that runs through all the results: 2025 is the
year with the highest ozone in the series, with 44.0\,\% station-day
exceedance versus 29.4\,\% in the fitting period, in all four seasons. It is
a year effect, not a calendar effect, and each technique revisits it where it
matters.

The five techniques are presented in the order in which each feeds the next:
PCA (\cref{sec:pca}) and FA (\cref{sec:analisis-factorial}) reduce the
15 predictors to three axes; the time-series characterization
(\cref{sec:serie-tiempo}) establishes the temporal dependence of those axes;
logistic regression and the discriminant (\cref{sec:clasificacion_tecs})
classify the station-day; and the autoregressive model
(\cref{sec:autorregresivo}) incorporates temporal dependence into the
classification.
